\section{Specialization to a common twist}\label{sec:common-twist}

The residue formula has independent twists. We must specialize them to one common twist without losing terms at colliding roots. This section proves that the scalar trace extends holomorphically to a generic common twist. It also identifies the polynomial relations preserved by that extension and gives a finite-jet representation for the specialized trace.

Keep the notation of Lemma~\ref{lem:residue}, and put
\[
 W=\prod_{i,b}V(t_i-b)\prod_{i<j}G(t_i-t_j).
\]
The nonzero parameter $K$ is the target common value of the independent twists $K_i$. The following four products depend only on the sites and the spectral argument $w$:
\begin{equation}\label{eq:common-twist-products}
\begin{aligned}
 d_{\rm full}(w)&=\prod_b[w-b+3][w-b+2],
 &w_1(w)&=\prod_b[3-w+b][2-w+b],\\
 w_2(w)&=\prod_b[w-b-1][w-b],
 &v(w)&=\prod_b[3-w+b][w-b].
\end{aligned}
\end{equation}
Here $d_{\rm full}$ denotes the denominator product, not an individual local transfer weight. With $Q(w)=\prod_i[w-t_i-1]$, define
\begin{equation}\label{eq:common-twist-L}
\begin{aligned}
 L(w)&=(-1)^N\left(Kw_1(w)\frac{Q(w+2)}{Q(w)}
             +K^{-1}w_2(w)\frac{Q(w-3)}{Q(w-1)}\right)\\
 &\quad+v(w)\frac{Q(w-2)Q(w+1)}{Q(w)Q(w-1)}.
\end{aligned}
\end{equation}
All site evaluations below mean \(w=b\); transfer formulas use
\(L/d_{\rm full}\). Let \(P\) be the Laurent polynomial part of \(L\)
in \(z=e^{2iw}\), consisting of the negative powers at zero and the
nonnegative powers at infinity. Its support lies in \([-N,N]\).
More precisely, every coefficient of \(P\) is a Laurent polynomial
on the entire root torus, with coefficients depending on the site
and twist parameters. To see this, put \(p_i=e^{2it_i}\).
After extracting the common root and spectral monomials, the two
root denominator polynomials associated with \(Q(w)\) and \(Q(w-1)\)
are
\[
 C_1(z)=\prod_i(z-q p_i),\qquad
 C_2(z)=\prod_i(z-q^2p_i),\qquad q=e^{2il}.
\]
Their leading coefficients are one and their constant coefficients
are \((-q)^N\prod_i p_i\) and \((-q^2)^N\prod_i p_i\), respectively.
These are units in \(\mathbb C[p_1^{\pm1},\ldots,p_N^{\pm1}]\).
Formal division at infinity uses the leading coefficients and
formal division at zero uses the constant coefficients. Both
therefore keep every coefficient in this Laurent ring, including
at collisions. The orders of the site products at the two ends
give the stated support bound.

In an ordinary separated configuration, put
$z_{i,s}=e^{2i(t_i+s)}=q^sp_i$. The partial fractions are normalized by
\[
 L(w)-P(w)=\sum_{i=1}^N\sum_{s=1}^2
                 \frac{r_{i,s}}{z-z_{i,s}}.
\]
Along $D_i=-K_i$ their coefficients are
\begin{equation}\label{eq:common-twist-pole-coefficients}
\begin{aligned}
 r_{i,1}&=(K-K_i)a_{i,1},
 &a_{i,1}&=(-1)^N i z_{i,1}[2l]\,w_1(t_i+1)
       \prod_{j\ne i}\frac{[t_i-t_j+2]}{[t_i-t_j]},\\
 r_{i,2}&=(K^{-1}-K_i^{-1})a_{i,2},
 &a_{i,2}&=(-1)^{N+1} i z_{i,2}[2l]\,w_2(t_i+2)
       \prod_{j\ne i}\frac{[t_i-t_j-2]}{[t_i-t_j]}.
\end{aligned}
\end{equation}
Indeed
\[
 \lim_{w\to t_i+s}\frac{z-z_{i,s}}{[w-t_i-s]}=i z_{i,s}.
\]
The self numerator is \([2l]\) for the first term of \(L\) at
$s=1$ and \(-[2l]\) for the second term at $s=2$, giving the
displayed factors. The first-to-third and second-to-third residue
ratios are respectively \(K/D_i\) and \(D_i/K\): the site ratios
give \((-1)^N\prod S^{\mp1}\), the other-root ratios give
\(\prod_{j\ne i}R^{\mp1}\), and the self ratios are one.
Combining the residues and substituting \(D_i=-K_i\) proves
\eqref{eq:common-twist-pole-coefficients}. In particular, the factors
\((-1)^N i z_{i,1}[2l]\) and \((-1)^{N+1}i z_{i,2}[2l]\)
are nonzero constants times root monomials. They remain units at
every finite collision.

\begin{samepage}
\begin{proposition}[Common-twist regularity]\label{prop:common-twist}
Fix sites avoiding one another's shift orbits. On generic independent-twist fibers define
\begin{equation}\label{eq:common-twist-trace}
 \tau(f)=\sum_{D_i=-K_i}\frac{Wf}{J},
 \qquad J=\det\bigl(\partial_{t_j}\log D_i\bigr).
\end{equation}
The sum is over the regular configurations of Lemma~\ref{lem:residue}. Then the following assertions hold near a generic diagonal twist $K_1=\cdots=K_N=K\ne0$. In particular, the genericity assumption excludes the toric-end resonances
\begin{equation}\label{eq:common-twist-resonances}
 (-K)^r=q^{-r^2}\quad\text{or}\quad(-K)^r=q^{r^2},
 \qquad 1\le r\le N,\qquad q=e^{2il}.
\end{equation}
\begin{enumerate}
\item For every regular Laurent polynomial in the root coordinates, $\tau(f)$ extends holomorphically across the diagonal.
\item On the diagonal, the continued trace of any product of site evaluations $L(b)$ equals the trace of the corresponding product of $P(b)$.
\item After clearing $Q(w)Q(w-1)(L-P)$ into a polynomial in $z$ by root and spectral monomials as necessary, the specialized trace annihilates every Laurent-polynomial multiple of each of its coefficients.
\end{enumerate}
\end{proposition}
\end{samepage}


\subsection{The marked forms and their collision orders}

\begin{proof}[Proof of Proposition~\ref{prop:common-twist}]
For $N=0$, the empty configuration is the single fiber point,
$W=J=1$, and $\tau(f)=f$ on $\mathbb C$. Also
$L=P=K+K^{-1}+1$, so all three assertions are immediate. Assume
$N\ge1$ for the rest of this proof.

Let $\pi=(-D_1,\ldots,-D_N)$ be the map on the compact model at
$F=\varnothing$ in Lemma~\ref{lem:residue}. We first locate the part of its boundary that
can meet a neighborhood of the diagonal. We then prove that the forms
whose traces occur in the proposition are holomorphic over that
neighborhood.

\paragraph{Toric ends and neutral rays.}
In multiplicative coordinates, as a ratio tends to zero,
$S\to q^{-1}$ and $R\to q$; as it tends to infinity,
$S\to q$ and $R\to q^{-1}$. Suppose that along solutions tending to
the diagonal exactly $r$ root coordinates tend to zero. Multiply their
equations. All internal pair factors cancel by $R(d)R(-d)=1$,
the $Nr$ site factors tend to $q^{-Nr}$, and the $r(N-r)$ pair
factors with the other labels tend to $q^{r(N-r)}$. Hence
\[
 \prod_{i:\,p_i\to0}(-K_i)\longrightarrow q^{-r^2},
 \qquad\text{so}\qquad (-K)^r=q^{-r^2}.
\]
If no root tends to zero but $r$ tend to infinity, the same calculation
gives $q^{Nr-r(N-r)}=q^{r^2}$ and hence $(-K)^r=q^{r^2}$.
This proves the necessary resonances
\eqref{eq:common-twist-resonances}. Compactness now shows that, after
shrinking a neighborhood $U$ of a generic diagonal point, $\pi^{-1}(U)$
meets no toric end.

This exclusion is uniform for the punctured site families used in
Section~\ref{sec:repeated-moments}. Suppose the multiplicative site
coordinates stay in a fixed compact annulus as $\epsilon\to0$,
and their shift orbits are distinct for $0<|\epsilon|<\delta$.
For some nonempty label set $J$, the same subset-product calculation
forces
\[
 \prod_{i\in J}(-K_i)=q^{-|J|^2}
 \quad\text{or}\quad
 \prod_{i\in J}(-K_i)=q^{|J|^2}
\]
at a toric end. The limits of $S$ and $R$ used above are uniform
when the site coordinates stay in that annulus. A generic diagonal
avoiding \eqref{eq:common-twist-resonances} has one neighborhood
$U$ avoiding all these finitely many subset-product hypersurfaces,
independently of sufficiently small $\epsilon$. Hence no toric end
meets the fibers over this $U$ for any $0<|\epsilon|<\delta$.
For each such $\epsilon$, the sites still have distinct shift
orbits, so the neutral-ray argument below applies on the same
twist neighborhood. This statement concerns the punctured family;
the collision at $\epsilon=0$ requires the separate two-site order
calculation in Section~\ref{sec:repeated-moments}.

Use the ray notation $\lambda_C,e_i$ from
\eqref{eq:residue-level-order}. If a boundary ray divisor meets
$\pi^{-1}(U)$, then
\begin{equation}\label{eq:common-twist-neutrality}
 e_i=0\qquad\text{for every label }i.
\end{equation}
Indeed a positive or negative $e_i$ makes the regular map $D_i$
identically zero or infinity on that divisor, whereas every coordinate
of $U$ is finite and nonzero. Let $A$ be an outermost positive-weight
cluster anchored at a site $b$, with the site included as a label of
$A$, and write $\mathrm{free}(A)$ for its free labels. For any anchored $C\subseteq A$, put
$a_C=m_{-1}(C)$ and $b_C'=m_1(C)$, with offsets relative to $b$.
Summing \eqref{eq:common-twist-neutrality} over the free labels of $A$
cancels every internal pair charge and gives
\begin{equation}\label{eq:common-twist-outer-neutrality}
 0=\sum_{i\in\mathrm{free}(A)}e_i
   =\sum_{\substack{C\subseteq A\\ C\ {\rm anchored\ at}\ b}}
        \lambda_C(a_C-b_C').
\end{equation}
The equality is a weighted sum over all anchored depths, not an
equality at each depth. It follows in particular that the total order
of $\prod_i S(t_i-b)^{-1}$ on this ray is zero: its order is the
negative of the right side, summed over the outer clusters at $b$.

\paragraph{The forms whose traces occur.}
Put $z_b=e^{2ib}$. At an ordinary configuration, evaluating
\eqref{eq:common-twist-L} at a site gives
\[
 L(b)=(-1)^N K w_1(b)\prod_i S(t_i-b)^{-1},
\]
because $w_2(b)=v(b)=0$. For a multiset $\mathcal E$ of unreplaced
evaluations, put
$\mathcal S_{\mathcal E}=\prod_{b\in\mathcal E}\prod_iS(t_i-b)^{-1}$.
Up to regular factors depending only on parameters, the three families
of top forms needed below are
\begin{align*}
 \eta_f&=Wf\bigwedge_i dt_i,\\
 \eta^{\rm ev}_{i,s;b,f}
   &=Wf\,\frac{a_{i,s}}{z_b-z_{i,s}}\,
          \mathcal S_{\mathcal E}\bigwedge_i dt_i,\\
 \eta^{\rm coeff}_{i,s;f}
   &=Wf\,a_{i,s}\bigwedge_i dt_i ,
\end{align*}
where $f$ is a regular Laurent polynomial in the root coordinates.
The second family arises by telescoping a product of $L(b)$ and $P(b)$
one factor at a time and extracting $K-K_i$ or
$K^{-1}-K_i^{-1}$ from
\eqref{eq:common-twist-pole-coefficients}. For the third family,
clearing $Q(w)Q(w-1)(L-P)$ cancels the simple-fraction denominator;
its remaining coefficients multiply $a_{i,s}$ by regular Laurent
polynomials. Thus holomorphic traces for these three families imply
the three assertions: every marked trace is multiplied by a twist
difference vanishing on the diagonal.

\paragraph{Orders of the marked forms.}
At a collision depth write $m_k$ and $B=B(m)$ as in
\eqref{eq:residue-B}, and put $a=m_{-1}$ and $b'=m_1$ when it is
anchored. The bare form $W\bigwedge_i dt_i$ has logarithmic order
\begin{equation}\label{eq:common-twist-bare-orders}
\begin{array}{c|c}
\text{cluster}&\text{order}\\ \hline
\text{unanchored}&B-1\\
\text{anchored}&B-a-b'.
\end{array}
\end{equation}
This is the pair order and discrepancy from the residue proof, with
the site order $-\delta_{k,-1}-\delta_{k,1}$ of $V$.
If the cluster contains the marked label at offset $k$, the root
quotients in $a_{i,1}$ and $a_{i,2}$ add respectively
\begin{equation}\label{eq:common-twist-mark-orders}
 m_{k+2}-m_k+1,\qquad m_{k-2}-m_k+1.
\end{equation}
The $+1$ removes the marked label itself from the denominator count.
At an anchored depth, $w_1(t_i+1)$ adds one at each of $k=1,2$,
and $w_2(t_i+2)$ adds one at each of $k=-2,-1$. An evaluation kernel
$(z_b-z_{i,s})^{-1}$ subtracts one only at its matching site and
$k=-s$. For $s=2$ this is canceled by the numerator just listed.
The only unmatched negative kernel order is the first pole at its
matching site with $k=-1$.

Orders add with the depth weights. Regular Laurent factors have no
negative order here because toric ends have been excluded.
Equation~\eqref{eq:common-twist-outer-neutrality} removes the total
order of every factor $\mathcal S_{\mathcal E}$; it does not remove
those factors depth by depth. We will show that the remaining total
logarithmic order is strictly positive on every boundary ray meeting
$\pi^{-1}(U)$.

\paragraph{Bounds for the offset counts.}
Let $M=\max_km_k$. The following bounds include the small and zero
cases needed in this proof and in the two-site reduction:
\begin{equation}\label{eq:common-twist-B-bounds}
 B\ge M,\qquad B\ge2M-1,\qquad
 B\ge2M\ \text{if }M\ge3.
\end{equation}
For $M>0$, equality $B=M$ occurs only for a unit singleton or a
unit adjacent pair. Here is a proof. If some count is zero, the
variation estimate in the residue proof gives
$\Delta=\frac12\sum(m_k-m_{k+1})^2\ge M$. Equality for $B=M$
also forces all distance-two products to vanish. The contiguous,
unit-or-flat rise and fall in the equality case for $\Delta$ then
leave only the two stated supports. If every count is positive, the
two distance-two products incident with a peak already sum to at
least $2M$, so equality $B=M$ is impossible.

For the second and third bounds, the cases $M=0,1$ follow from
$B\ge M$. Suppose some count is zero and choose a peak at index $0$.
If $m_{-2}$ or $m_2$ is positive, its product with the peak adds
at least $M$ to $\Delta\ge M$. Otherwise the four gradients from
the peak to these two zero counts give
\[
 \Delta\ge\frac12\bigl((M-m_{-1})^2+m_{-1}^2
                       +(M-m_1)^2+m_1^2\bigr)\ge\frac{M^2}{2}.
\]
For $M\ge4$ this is at least $2M$. For $M=3$, equality in the
integral lower bound $5$ forces both first neighbors to be positive,
so their distance-two product adds at least one; otherwise the
gradient contribution is already larger. For $M=2$, the only way
those four gradients contribute just $2$ is $m_{-1}=m_1=1$,
again adding their product; otherwise they contribute at least $3$.
This proves \eqref{eq:common-twist-B-bounds}.

\paragraph{Nonnegative depth bounds and their outer equality cases.}
For an unanchored marked depth, combining
\eqref{eq:common-twist-bare-orders} and
\eqref{eq:common-twist-mark-orders} gives
$B-m_k+m_{k\pm2}\ge0$. The unmarked bound is $B-1\ge0$.
Since an unanchored cluster has at least two labels, equality in
either bound requires a unit adjacent pair.

For an anchored unmarked depth, $B-a-b'\ge0$ by the stage-$0$
bound in the residue proof. Apart from the matching first kernel,
discarding the nonnegative numerator orders gives the marked lower
bound
\begin{equation}\label{eq:common-twist-anchored-bound}
 B-a-b'-m_k+1.
\end{equation}
If at most one of $a,b'$ is positive, this is nonnegative by
$a+b'\le M$, $m_k\le M$ and $B\ge2M-1$. If both are positive,
\[
 B\ge M+ab',\qquad
 B-a-b'-m_k+1\ge (M-m_k)+(a-1)(b'-1)\ge0.
\]
The first inequality follows from $\Delta\ge M$ and the
distance-two product $ab'$ when some count is zero. When all counts
are positive it is strict, because another product incident with
a peak supplies $M$ in addition to $ab'$ and further products remain.
Equality requires no distance-two products other than $ab'$ and a
contiguous monotone support with steps of size at most one. Hence
the support is exactly $\{-1,0,1\}$, with
$m_{-1}=m_1=1$ and $1\le m_0\le2$.

At an outer anchored cluster, \eqref{eq:common-twist-outer-neutrality}
rules out exactly one of $a,b'$ being positive: all nested counts
at the absent end would still be zero, so the weighted sum would
have the strict sign of the other end. If both vanish, the
unmarked and marked bounds are strictly positive by $B\ge M$.
Thus any outer equality in the bounds just considered has support
$\{-1,0,1\}$ with unit end counts. The unmarked equality is the
subcase $m_0=1$.

\paragraph{The matching first kernel.}
It remains to handle an evaluation mark with $s=1$ at its own site
and offset $-1$. The anchored clusters inside an outer anchored
cluster form a chain, because they all contain the site. Along the
initial part of that chain containing the marked label, the exact
bound is
\[
 (B-a-b')+(b'-a+1)-1=B-2a.
\]
After the marked label leaves the anchored chain, the bound is
$B-a-b'$. If $b'>0$ at every marked depth, then
$B\ge M+ab'\ge2a$ makes all these bounds nonnegative; equality
at the outer depth again has the support and unit ends above.

Otherwise some marked anchored depth has $b'=0$. Every deeper
anchored depth then has $b'=0$ by nesting, including the depths
after the mark has left. Add the zero quantity in
\eqref{eq:common-twist-outer-neutrality} to the total anchored
order for this outer cluster. The bounds become $B-a-b'$ while
the mark is anchored and $B-2b'=B$ afterwards. They are all
nonnegative. In a total equality, the outer bound is $B-a-b'=0$;
outer neutrality and the marked label at $-1$ again force the
same support with unit ends.

\paragraph{Strict order on a neutral ray.}
Suppose the total logarithmic order of one of the three forms were
zero. The preceding nonnegative bounds would have to be equalities
at every positive depth. If an outer cluster were unanchored, its
unit adjacent pair would have pair charges $+1,-1$ and no strict
subcluster, contradicting \eqref{eq:common-twist-neutrality}.
Every remaining outer equality has support $\{-1,0,1\}$, unit
counts at $-1,1$, and positive $m_0$. For its unique label at
$-1$, the charge from this outer depth is
\[
 \underbrace{m_{-3}+m_0-m_{-2}-m_1}_{\text{pair charge}}
 +\underbrace{1}_{\text{site charge}}=m_0>0.
\]
At every deeper anchored depth containing that label the charge is
$1+m_0-m_1\ge0$, since the count at $1$ is at most one. A deeper
unanchored equality containing it must be an adjacent unit pair
within the same support, necessarily using an offset-$0$ label,
and contributes $+1$. Thus the total $e_i$ is positive, again
contradicting neutrality. Every boundary ray meeting $\pi^{-1}(U)$
therefore has strictly positive logarithmic order for each form.
Since these orders are integral, the pulled-back forms are
holomorphic there, including at intersections of boundary divisors.

\subsection{From holomorphic forms to scalar traces}

It remains to justify that dividing by the Jacobian and summing over
the fibers preserves holomorphy. The restriction
$\pi:\pi^{-1}(U)\to U$ is proper. Let $\eta$ be any of the
holomorphic top forms just obtained, and put
$\Omega=\bigwedge_i d\log K_i$, a nonvanishing volume form on $U$.
On regular fibers define
\[
 T_\eta(\mathbf K)=
 \sum_{x\in\pi^{-1}(\mathbf K)}
       \frac{\eta}{\pi^*\Omega}(x).
\]
For $\eta=\eta_f$, the denominator in root coordinates is $J$, so this
is exactly $\tau(f)$.

First, $T_\eta$ is a rational function on the generic target.
If the target coordinates $-D_i$ are algebraically dependent, the
generic fiber is empty and this trace is identically zero, so there is
nothing further to prove. Otherwise
$\mathbb C(p_1,\ldots,p_N)$ is a finite algebraic extension of
$\mathbb C(-D_1,\ldots,-D_N)$, by equality of transcendence degrees.
It is separable in characteristic zero. Adjoin the $p_i$ successively:
outside the zeros of the minimal-polynomial discriminants and the
relevant denominators, their distinct roots give precisely the
generic inverse branches. A nonzero input denominator can be
excluded by the nonzero norm of that field element on the target.
Summing the rational function $\eta/\pi^*\Omega$ over the branches
is symmetric at each field step, and is therefore rational. This
also gives a finite generic degree $d_\pi$ bounding the number of
branches.

Now take a relatively compact coordinate polydisc $V\Subset U$.
On its regular values, Cauchy--Schwarz and change of variables give
\begin{equation}\label{eq:common-twist-L2}
 \int_{V_{\rm reg}}|T_\eta|^2\,dV_\Omega
 \ \le\ d_\pi\int_{\pi^{-1}(V)}dV_\eta
 \ <\ \infty,
\end{equation}
where $dV_\Omega=i^{N^2}\Omega\wedge\overline\Omega$ and
$dV_\eta=i^{N^2}\eta\wedge\overline\eta$. Indeed
$|\sum_{\ell=1}^{d_\pi}a_\ell|^2\le
d_\pi\sum_\ell|a_\ell|^2$ for the branch values
$a_\ell=\eta/\pi^*\Omega$, and the squared complex Jacobian in
change of variables cancels their denominators. The last integral
is finite because $\eta$ is holomorphic on the compact preimage
of $\overline V$. Positive-dimensional exceptional fibers cause
no problem, since the estimate uses only the regular fibers and
the generic degree.

A rational function with a pole along a divisor is not locally
square-integrable: at a generic smooth point of that divisor its
form is $z_1^{-m}$ times a holomorphic unit, $m\ge1$, and the
integral in the transverse complex coordinate diverges. Thus
\eqref{eq:common-twist-L2} excludes every polar divisor of
$T_\eta$; cancellation of the irreducible factors of a local denominator
then shows that it is holomorphic on $U$. The same argument works
jointly with sites in a small neighborhood avoiding one another's
shift orbits. The compactification is then made relatively in the
site coordinates; wedging $\eta$ with their differentials and using
the proper map $(\mathbf p,\mathbf b)\mapsto(-\mathbf D,\mathbf b)$
gives equal dimensions and the same estimate.
In the punctured family just considered, these local proper-trace
arguments cover the common product neighborhood
$U\times\{0<|\epsilon|<\delta\}$. Consequently, the trace of a Laurent
insertion whose coefficients are jointly holomorphic in the sites,
twists and any row parameters is a finite holomorphic combination
of the rational monomial traces, and is holomorphic there. When
those coefficients are also rational in the parameter coordinates,
the trace is rational and has no polar divisor in that product
neighborhood. Any inverted scalar row denominators must stay
regular and nonzero. This supplies the off-$\epsilon=0$ input to
the later two-site proof, without asserting an order at the
collision itself.

Apply this result to the three families of forms. It gives the
holomorphic continuation of $\tau(f)$. Telescoping a product of
site evaluations expresses its difference from the corresponding
product of $P(b)$ as a sum of the holomorphic marked traces
multiplied by $K-K_i$ or $K^{-1}-K_i^{-1}$. Their restrictions to
the diagonal vanish. After clearing the spectral denominators,
the same argument for every Laurent-polynomial multiple of a
coefficient gives its annihilation by the specialized trace.
This proves all three assertions.
\end{proof}

\subsection{A finite-jet representation of the specialized trace}

The collision limit need not be a sum of ordinary evaluations. The next corollary records the additional derivatives that can remain, and the relations satisfied at every point that contributes. It is used at generic $K$ satisfying the preceding hypotheses. Section~\ref{sec:physical-amplitude} reaches the physical values $K=i$ and $K=1$ by analytic continuation of the physical branch and its coefficient, after using this corollary at generic nearby twists. Write $\mathbf p=(p_1,\ldots,p_N)$, where $p_i=e^{2it_i}$, and let
\[
 \mathcal A=\mathbb C[p_1^{\pm1},\ldots,p_N^{\pm1}]
\]
be the Laurent polynomials regular on the entire root torus.

\begin{corollary}[Finite jets at a common twist]\label{cor:common-twist-jets}
Fix the sites and a generic common twist as in Proposition~\ref{prop:common-twist}, and denote the specialized trace by $\tau_0$. There are finitely many points $\mathbf p^{(\nu)}\in(\mathbb C^*)^N$, integers $m_\nu\ge0$, and constants $c_{\nu,\alpha}$ such that, for every $f\in\mathcal A$,
\begin{equation}\label{eq:common-twist-jets}
 \tau_0(f)=\sum_\nu\ \sum_{|\alpha|\le m_\nu}
 c_{\nu,\alpha}\,\partial_{\mathbf p}^{\alpha}f(\mathbf p^{(\nu)}).
\end{equation}
Here $\alpha\in\mathbb Z_{\ge0}^N$, $|\alpha|=\sum_i\alpha_i$, and repeated points are combined and zero local functionals omitted. Every remaining point satisfies the coefficients of the cleared relation $Q(w)Q(w-1)(L-P)=0$; consequently $L=P$ there as a rational identity in the spectral argument. Such a point may have colliding roots. At a regular transverse solution in the original configuration locus, the contribution is the ordinary trace term $Wf/J$, with $J$ as in \eqref{eq:common-twist-trace}; its weight is nonzero whenever $W\ne0$.
\end{corollary}

\begin{proof}
For $N=0$, the root torus is the single empty tuple,
$\mathcal A=\mathbb C$, and $\tau_0(f)=f$. The representation is
one order-zero jet with coefficient one, and the cleared relation
vanishes because $L=P$. Assume $N\ge1$ below. If the generic
fiber is empty, the trace is zero and the empty sum gives the
representation, so we may also assume that the map is dominant.

Choose a generic affine line through the fixed diagonal point, with
parameter $\epsilon=0$ there. After shrinking the disk, its nonzero
points avoid the exceptional algebraic locus of the finite generic
fibers. The finitely many inverse branches have finite permutation
monodromy, so after a power substitution $\epsilon=s^m$ they are
single-valued on the punctured $s$-disk. Properness and the exclusion
of toric ends keep every root coordinate bounded above and away from
zero. Each branch therefore extends holomorphically to $s=0$, with
a limit in $(\mathbb C^*)^N$; collisions are allowed.

The branch weights $W/J$ are rational functions of the roots and
parameters. On each extended branch they are meromorphic in $s$,
with some finite pole order. Choose one integer $M$ bounding those
orders for all branches. It is independent of $f$. Since the trace
is holomorphic at $\epsilon=0$, $\tau_0(f)$ is the constant term in
$s$ of the sum of the weighted branch values $f(\mathbf p(s))$.
The Taylor expansion at a branch limit contributes to this constant
term only through derivatives of order at most $M$, because
$\mathbf p(s)-\mathbf p(0)=O(s)$. This gives
\eqref{eq:common-twist-jets}; combine equal limits and discard zero
local functionals.

Let $c(\mathbf p)$ be any coefficient of the cleared relation.
Proposition~\ref{prop:common-twist} gives $\tau_0(cf)=0$ for every
$f\in\mathcal A$. The powers of the distinct maximal ideals at the
finitely many support points are pairwise comaximal, so Laurent
polynomials can prescribe their finite jets independently. If
$c(\mathbf p^{(\nu)})\ne0$, multiplication by $c$ is invertible in
the finite jet algebra at that point. We could then prescribe an
arbitrary jet of $cf$ there and zero jets at the other points,
forcing the local functional at $\nu$ to vanish. This contradicts
its retention in the sum. Thus every cleared coefficient vanishes
at every support point. Since $Q(w)Q(w-1)$ is a nonzero rational
function of the spectral argument for torus roots, the cleared
identity gives $L=P$ as a rational identity there.

Finally, at a regular transverse solution in the original
configuration locus the implicit function theorem gives a unique
ordinary local branch. Its weight is holomorphic with value $W/J$,
so its constant term is precisely $(W/J)f$ at that point, with no
derivatives. It is nonzero when $W\ne0$, as asserted.
\end{proof}
